\subsection{图形与对应关系}

定义 1. 若 G 的每个元素都是有序对，即若关系

\[
(\forall z) (z \in G \Longrightarrow (z \text {   是有序对 }))
\]

为真。

若 G 是一个图形，则关系 \((x, y) \in \mathbf{G}\) 表示为“在 G 下 y 对应于 x”。

让 \(R \mid x, y \mid\) 是一个关系，其中 \(x\) 并且 \(y\) 是不同的字母。设 \(G\) 是一个字母，不同于 \(x\) 并且 \(y\) 均不同，且不出现在 \(R\) 中的字母。如果关系

\[
(\exists G) (G \text {   是图   且   } (\forall x) (\forall y) (R \Longleftrightarrow ((x, y) \in G)))
\]

为真时，关系R被称为具有一个图（关于字母x和y）。根据外延公理，图G是唯一的，并被称为R关于x和y的图。

让 \(\pmb{Z}\) 是一个字母，不同于 \(\pmb{x}\) 并且 \(\pmb{y}\) 均不同，且不出现在 \(\pmb{R}\) 中的字母。如果关系

\[
(\exists Z) (\forall x) (\forall y) (R \Longrightarrow ((x, y) \in Z))
\]

成立，则 \(R\) 具有一个图；我们可以取这个图为有序对的集合 \(z\) 使得 \(z \in Z\) 并且 \(R \{ \text{pr}_1 z, \text{pr}_2 z \}\) ( \(z\) 是一个不同于 \(x, y, Z\) 的字母，它不出现在 \(R\) 中）。如果我们知道一个项，则此条件被满足 \(T\) ，它不包含 \(x\) 要么 \(y\) 之前，使得 \(R \Longrightarrow ((x, y) \in T)\) 为真。

命题1. 设G是一个图。存在唯一一个集合A和唯一一个集合B，它们具有以下性质：

\begin{enumerate}
\def\labelenumi{(\arabic{enumi})}
\item
  该关系 \((\exists y)((x, y) \in G)\) 等价于 \(x \in A\) ;
\item
  该关系 \((\exists x)((x, y) \in \mathbf{G})\) 等价于 \(y \in \mathbf{B}\) .
\end{enumerate}

因为只需取A（相应地，B）为所有形如 \(pr_{1}z\) （分别） \(pr_{2}z\) 的对象所组成的集合，其中 \(z \in G\) （§1，第6节）。确切地说，

\[
\mathrm{A} = \mathcal {E} _ {x} ((\exists y) ((x, y) \in \mathrm{G}))) \quad \text { 且 } \quad \mathrm{B} = \mathcal {E} _ {y} ((\exists x) ((x, y) \in \mathrm{G}));
\]

这些集合分别称为图G的第一投影和第二投影，或G的定义域和值域；它们记作 \(pr_{1}\langle G\rangle\) 并且 \(pr_{2}\langle G\rangle\) （或记作 \(pr_{1}G\) 并且 \(pr_{2}G\) ，如果没有混淆的风险）。立即可以验证 \(G\subset(\mathrm{pr}_{1}\mathrm{G})\times(\mathrm{pr}_{2}\mathrm{G})\) ；因此，每一组有序对都是某个乘积的子集，反之亦然。如果这两个集合中的某一个 \(pr_{1}G\) , \(pr_{2}G\) 为空，我们有 \(G=\varnothing\) （§2，命题2）。

备注。该关系 \(x = y\) 没有图，因为如果图存在的话，其第一投影将是所有对象的集合（参见§1，第7号，注记）。

\textbf{定义2。} 集合A与集合B之间的一个对应是一个三元组

\[
\Gamma = (\mathrm{G}, \mathrm{A}, \mathrm{B}),
\]

其中G是一个图，使得 \(pr_{1}G \subset A\) 并且 \(pr_{2}G \subset B\) 。称 G 为 \(\Gamma\) 的图，A 是源，B 是 \(\Gamma\) 的目标.

如果 \((x, y) \in G\) ，我们说“y 在对应关系 \(\Gamma\) 中对应于 x”。对于每个 \(x \in pr_{1}G\) ，该对应 \(\Gamma\) 被称为在x处有定义，且 \(pr_{1}G\) 被称为 \(\Gamma\) 的定义域；每个 \(y \in pr_{2}G\) 称为 \(\Gamma\) 所取的一个值，以及 \(pr_{2}G\) 称为 \(\Gamma\) 的值域.

如果 \(R\{x, y\}\) 是一个具有图像的关系 \(G\) （就字母而言） \(x, y\) )，且如果 \(A, B\) 是两个集合，使得 \(\operatorname{pr}_1 G \subset A\) 并且 \(\operatorname{pr}_2 G \subset B\) ，我们说R是A中元素与B中元素之间的一种关系（关于字母x，y）。这种对应关系 \(\Gamma = (G, A, B)\) 被称为由关系R定义的A与B之间的对应关系（关于x和y）。

让 \(\mathbf{G}\) 设为一个图，并令 \(\mathbf{X}\) 是一个集合。关系

\[
x \in \mathrm{X} \text {   且   } (x, y) \in \mathrm{G}
\]

蕴含 \((x, y) \in G\) ，因而有一个图 \(G'\) 。 \(G'\) 的第二投影由在 G 下与 X 中对象相对应的所有对象组成。

定义 3. 设 G 是一个图，X 是一个集合。所有在 G 下与 X 的元素相对应的对象的集合称为 X 在 G 下的像，记为 G⟨X⟩ 或 G(X)。

让 \(\Gamma = (\mathbf{G},\mathbf{A},\mathbf{B})\) 是一个对应关系，且设 \(\mathbf{X}\) 为 \(\mathbf{A}\) 。集合 \(\mathbf{G}\langle \mathbf{X}\rangle\) 也记为 \(\Gamma \langle \mathbf{X}\rangle\) 要么 \(\Gamma (\mathbf{X})\) 并且被称为 \(\mathbf{X}\) 在……之下 \(\Gamma\) .

注记

\begin{enumerate}
\def\labelenumi{(\arabic{enumi})}
\item
  精确地说， \(\mathbf{G}\langle \mathbf{X}\rangle\) 表示集合 \(\mathcal{E}_{\gamma}((\exists x)(x\in \mathbf{X}\) 并且 \((x,y)\in \mathbf{G}))\) 。从现在起，我们通常不再将定义翻译成形式语言。
\item
  记号 G(X) 和 \(\Gamma(X)\) 偶尔会与后面引入的记号产生混淆（参见第 4 号，定义 9 之后的注记）。
\end{enumerate}

设 G 是一个图。由于关系 \((x, y) \in G\) 蕴含 \(y \in pr_{2}G\) ，对每个集合 X 都有 \(G\langle X \rangle \subset pr_{2}G\) 。由于 \((x, y) \in G\) 蕴含 \(x \in pr_{1}G\) ，故 \(G\langle pr_{1}G \rangle = pr_{2}G\) 。我们有 \(G\langle \phi \rangle = \phi\) ，因为 \(x \notin \phi\) 是一个定理。如果 \(X \subset pr_{1}G\) 且 \(X \neq \phi\) ，则 \(G\langle X \rangle \neq \phi\) .

命题2. 设 G 是一个图，X, Y 是两个集合；则关系 \(X \subset Y\) 蕴含 \(G\langle X\rangle \subset G\langle Y\rangle\) .

这是定义以及 C50（§1，第4节）的直接推论。

推论. 如果 \(\mathbf{A} \supset \mathrm{pr}_1\mathbf{G}\) ，我们有 \(\mathbf{G}\langle \mathbf{A} \rangle = \mathrm{pr}_2\mathbf{G}\) .

定义4. 设 G 是一个图且 x 是一个对象。集合 G\textless\{x\}\textgreater（有时为了语言上的简便，也记作G(x)）称为G在x处的截面。

由C43（第一章，第5节，第1号）可直接得出，该关系 \(y \in G \langle \{x\} \rangle\) 等价于 \((x, y) \in G\) 。如果 \(G\) 并且 \(G'\) 是两个图，关系 \(G \subset G'\) 因此等价于

\[
(\forall x) (\mathrm{G} \langle \{x \} \rangle \subset \mathrm{G} ^ {\prime} \langle \{x \} \rangle).
\]

如果 \(\Gamma = (G, A, B)\) 是...之间的对应关系 \(A\) 并且 \(B\) ，那么对于每一个 \(x \in A\) 的截面 \(G\) 在 \(x\) 也称为...的截面 \(\Gamma\) 在 \(x\) 并且用 \(\Gamma \langle \{x\} \rangle\) （或 \(\Gamma(x)\) ).

